\chapter{Paradoxos de Tamanho e Contagem}
\label{chap:apendiceTamanhoContagem}

Os paradoxos reunidos neste apêndice formam uma narrativa histórica
que parte da Grécia Antiga e culmina na teoria cantoriana, cada um deles revelou uma insuficiência no entendimento vigente do infinito e motivou avanços fundamentais na matemática.

\section{Os Paradoxos de Zenão (c. 450 a.C.)}

Apresentar Aquiles e a tartaruga, a flecha em repouso, e a dicotomia. Mostrar que a inquietação com somas infinitas e divisibilidade infinita é anterior à matemática formal em mais de dois milênios.


\section{O Paradoxo de Galileu (1638)}

Apresentar a bijeção \(n \to n^2\) entre os naturais e os quadrados perfeitos. Contextualizar como o primeiro registro moderno de estranheza com bijeções em conjuntos infinitos. Conectar com a Definição de conjunto infinito do Capítulo 3.


\section{O Hotel de Hilbert (1924)}

Apresentar o experimento mental: hotel com infinitos quartos, todos ocupados, que ainda assim acomoda novos hóspedes inclusive infinitos deles. Conectar com a aritmética cardinal do Capítulo 6 (\(ℵ_0 + ℵ_0 = ℵ_0\), etc).


\section{O Paradoxo de Cantor (1899)}

Apresentar a contradição: o conjunto de todos os conjuntos teria que ter cardinalidade maior que si mesmo via Teorema de Cantor. Conectar com o Teorema 1 (não existe conjunto de todos os conjuntos) já demonstrado no Capítulo 2.


\section{O Paradoxo de Burali-Forti (1897)}

Apresentar a contradição sobre o conjunto de todos os ordinais: tal conjunto seria ele mesmo um ordinal, e seu sucessor seria simultaneamente maior e não maior que ele. Situar como o primeiro paradoxo formal da teoria dos conjuntos, anterior inclusive ao de Russell. Conectar com os Números Ordinais do Capítulo 7.
